%& -job-name=newfilenameialwayswanted
\documentclass{beamer}
% (fold)
\graphicspath{{./Images/}} % the folder in which images are stored for this project
\usepackage{tikz,graphicx,amsmath,hyperref,tcolorbox,fancyvrb,xcolor}
\hypersetup{colorlinks=true} % make sure our hyperlinks are coloured for visibility
\definecolor{DBlack}{RGB}{35,31,32}
\definecolor{DPurple}{RGB}{126,49,123}
\definecolor{DBlue}{RGB}{0,99,136}

\tcbuselibrary{skins,breakable}
\newenvironment{BGVerbatim}
 {\VerbatimEnvironment
  \begin{tcolorbox}[
    breakable,
    colback=lightgray,
	boxsep=3mm
  ]%
  \begin{Verbatim}}
 {\end{Verbatim}\end{tcolorbox}}

\usecolortheme[RGB={83,13,88}]{structure}
% This is the theme that we are modifying
\usetheme{Boadilla}

% Change the font colour to white on the title page to stand out against the purple
\setbeamercolor{title}{fg=white}
\setbeamerfont{title}{family=\rmfamily}
\setbeamercolor{author}{fg=white}
\setbeamerfont{author}{series=\bfseries}
\setbeamercolor{institute}{fg=white}
\setbeamercolor{date}{fg=white}
\setbeamerfont{date}{series=\bfseries}

% Remove the default ball bullet points in the toc, but keep the numbering
\setbeamertemplate{sections/subsections in toc}[sections numbered]
% Change the item icons to be circles rather than balls
\setbeamertemplate{itemize items}[circle]
% Define blocks to be rounded rectangles with shadows
\setbeamertemplate{blocks}[rounded][shadow=true]

% Turn the @ symbol into character class other
\makeatother
% Modify the footline to get a two part footline with slide numbers
\setbeamertemplate{footline}
{
  \leavevmode%
  \hbox{%
  % Create a box with width of slide, containing current slide and total slides
  \begin{beamercolorbox}[wd=\paperwidth,ht=2.25ex,dp=1ex,center]{section in head/foot}%
    \insertframenumber{} / \inserttotalframenumber\hspace*{1ex}
  \end{beamercolorbox}}%
  \vskip0pt%
}
% Turn the @ symbol back to a letter
\makeatletter

% Remove the default navigation symbols from the bottom of the slides
\setbeamertemplate{navigation symbols}{}
% (end)
\author{Sam Fearn\\{\small 
(s.m.fearn@durham.ac.uk)}}
%\institute[Durham]{Durham University}
\title{Building a website in 10 minutes with Jekyll}
\date{November 21\textsuperscript{st}, 2018}

\begin{document}
% (fold)
% --- Start Custom Commands --- %
\renewcommand{\d}[2][]{\ensuremath{\operatorname{d}^{#1}\!{#2}}}
% --- End Custom Commands --- %

%--- the titlepage frame -------------------------%
{
% Draw a purple box for the background of the title slide and add the durham logo at the bottom
\setbeamertemplate{background canvas}{
\begin{tikzpicture}
    \clip (0,0) rectangle (\paperwidth,\paperheight);
    \fill[color=DPurple] (0,0) rectangle (\paperwidth,\paperheight);
\end{tikzpicture}
}

\begin{frame}[plain]
\maketitle
\end{frame}
}

%---Complete Outline--------------------%
\begin{frame}
       \frametitle{Outline}
       \tableofcontents
\end{frame}
% (end)
\Large
\section{Motivation} % (fold)
\label{sec:motivation}

\begin{frame}{My Motivation}
	Why I built a website:
	
	\begin{itemize}
		\item<2-> Needed a place to put the details of my project for students to read.
		\item<3-> Wanted a place to be able to make notes to myself that others might find interesting.
		\item<4-> Didn't want anything complicated.
	\end{itemize}
\end{frame}

\begin{frame}{Your Motivation?}
	\Large
	Why \textbf{you} might want a website:
	
	\begin{itemize}
		\item<2-> Hold information about your teaching and projects?
		\item<3-> Host a blog about things you're working on / thinking about?
		\item<4-> Host your CV and/or publication details?
		\item<5-> \ldots ?
	\end{itemize}
\end{frame}

\begin{frame}{Demo 1}
	I wrote two posts discussing most of what I'm going to cover here on my site.
	
	\medskip
	\begin{center}
		\url{http://www.maths.dur.ac.uk/~sxwc62/}
	\end{center}
	
	\medskip
	\pause
    \begin{columns}
	\begin{column}{0.4\textwidth}
		This is a \emph{static} site built with Jekyll. I'll explain what this means later.
	\end{column}
	\begin{column}{0.5\textwidth}
		\centering
		\includegraphics[width=0.9\textwidth]{jekyllvial}
	\end{column}
    \end{columns}
	
\end{frame}

% section motivation (end)

\section{Building Blocks}
\label{sec:bblocks}

\subsection{The big picture} % (fold)
\label{sub:the_big_picture}
\begin{frame}{A Modern Complicated Website}
	\centering
	\includegraphics[width=\textwidth]{clientserver}
\end{frame}

\begin{frame}{A Modern Simple Website}
	\centering
	\includegraphics[width=\textwidth]{clientnoserver}
\end{frame}

\begin{frame}{Client-Side Technologies}
    \begin{columns}
	\begin{column}{0.6\textwidth}
		\centering
		\includegraphics[height=0.8\textheight]{clienttech}
	\end{column}
	\begin{column}{0.4\textwidth}
		\begin{itemize}
			\item<2->HTML - The Content
			\item<3->CSS - The Presentation
			\item<4->Javascript - (Locally) Dynamic Content
		\end{itemize}
	\end{column}
    \end{columns}
\end{frame}

% subsection the_big_picture (end)

\subsection{HTML} % (fold)
\label{sub:html}

\begin{frame}{HTML}
	\begin{itemize}
		\item<1-> HTML (or HyperText Markup Language) is the language we use to create the \textbf{content} and \textbf{structure} of our website.\pause This is separate to the \textbf{style}.
		\item<3-> Like a \LaTeX{} document (.tex), an html page is just a plain text file consisting of commands which are to be interpreted later. For \LaTeX{} the .tex file is interpreted by pdflatex (or other) to produce a pdf which we read.
		\item<4-> An HTML file (.html) is processed by a web browser in order to create visual document from the structured one.
	\end{itemize}
	
\end{frame}

\begin{frame}[fragile]
\normalsize
	We're all familiar with creating a basic \LaTeX{} document.
\begin{BGVerbatim}[numbers=left,numbersep=10pt]
\documentclass{article}
% We define an Author, Title and Date
\author{Sam Fearn}
\title{A Very Quick Introduction To \LaTeX{}}
\date{November 21\textsuperscript{st}, 2018}

\begin{document}
% Create a title from our Author, Title and Date
\maketitle
\section{Introduction}
Some introductory text goes here
\section{Content}
The main content goes here
\end{document}
\end{BGVerbatim}
\end{frame}

\begin{frame}[fragile]
\normalsize
	The equivalent \textbf{content} in HTML might look like:
\begin{BGVerbatim}[numbers=left,numbersep=10pt]
<!DOCTYPE html>
<html>
<body>

<h1>A Very Quick Introduction To HTML</h1>
<h3>Sam Fearn</h3>
<h3>November 21<sup>st</sup> 2018</h3>

<h2>Introduction</h2>
<p>Some introductory text goes here</p>
<h2>Content</h2>
<p>The main content goes here</p>

</body>
</html>
\end{BGVerbatim}	
\end{frame}

\begin{frame}[fragile]{HTML Tags}
	HTML contains many tags, including:
	\begin{itemize}
		\item<1-> Opening \verb|<head>| and closing \verb|</head>| tags, used to contain various types of metadata.
		\item<2-> The anchor tag \verb|<a href="link">|, used to add links to other web pages.
		\item<3-> The \verb|<img src="link">| tag, used to add images to pages.
		\item<4-> (Unordered) Lists can be added by using \verb|<ul> </ul>| to denote the list, and then \verb|<li> </li>| to add elements to the list.
	\end{itemize}
\end{frame}

% subsection html (end)

\subsection{Markdown} % (fold)
\label{sub:markdown}

\begin{frame}{Markdown}
    \begin{columns}
	\begin{column}{0.6\textwidth}
		\begin{itemize}
			\item<1-> Markdown (.md) is processed to create HTML.
			\item<2-> It covers a subset of the HTML tags, with a focus on the text content.
			\item<3-> Markdown is very easy to read in plain text form.
			\item<4-> Can include regular HTML.
			\item<5-> Markdown can be converted to \LaTeX{}.
		\end{itemize}
	\end{column}
	\begin{column}{0.4\textwidth}
		\centering
		\includegraphics[width=0.9\textwidth]{Markdown}
	\end{column}
    \end{columns}
\end{frame}

\begin{frame}[fragile]{Markdown}
\normalsize
	The same example as before, now in Markdown:
\begin{BGVerbatim}[numbers=left,numbersep=10pt]
# A Very Quick Introduction To HTML

### Sam Fearn

### July 20th 2018

## Introduction

Some introductory text goes here

## Content

The main content goes here
\end{BGVerbatim}
\end{frame}

\begin{frame}
	Jekyll uses markdown as the main filetype for producing the content of a website and processes these into HTML files during the Jekyll `build phase'.
	
	\centering
	\includegraphics[width=0.5\textwidth]{jekyllvial}
\end{frame}

% subsection markdown (end)

\subsection{CSS} % (fold)
\label{sub:css}

\begin{frame}{CSS}
    \begin{columns}
	\begin{column}{0.6\textwidth}
		\begin{itemize}
				\item Cascading Style Sheets (CSS) are used to define the style and formatting for a website.
				\item CSS defines the \textbf{style \& positioning} to apply to the \textbf{content \& structure} defined in HTML.
				\item We can save our CSS as a .css file and link to it in the head of our HTML page.
			\end{itemize}
	\end{column}
	\begin{column}{0.4\textwidth}
		\centering
		\href{http://blog.int3ractive.com/2012/04/tutorial-realistic-waterfall-with-css3.html}{\includegraphics[width=0.9\textwidth]{waterfall}}
	\end{column}
    \end{columns}
\end{frame}

\begin{frame}[fragile]{A Basic CSS Example}
\normalsize
	A very simple CSS file might contain:
\begin{BGVerbatim}[numbers=left,numbersep=10pt]
body {
    background-color: lightgray;
}

h1 {
    color: purple;
    text-align: center;
}

p {
    font-family: Arial, Helvetica, sans-serif;
    font-size: 18px;
}
\end{BGVerbatim}
\end{frame}

\begin{frame}{CSS Features}
	\begin{itemize}
		\item<1-> Although the previous example was very simple, CSS is very powerful.
		\item<2-> CSS is capable of defining layout and positioning commands.
		\item<3-> CSS can provide basic animation.
		\item<4-> We can also use CSS to apply different styles depending on the browser that is trying to view our page -- this lets us support mobile devices.
	\end{itemize}
\end{frame}

% subsection css (end)

\subsection{SASS} % (fold)
\label{sub:sass}

\begin{frame}{Sass}
    \begin{columns}
	\begin{column}{0.6\textwidth}
		\begin{itemize}
				\item<1-> Sass (.scss) is a CSS preprocessor -- just like Markdown for HTML.
				\item<2-> Enables features not part of standard CSS, such as functions, variables and includes
		\end{itemize}
	\end{column}
	\begin{column}{0.4\textwidth}
		\centering
		\includegraphics[width=0.9\textwidth]{sass}
	\end{column}
    \end{columns}	
\end{frame}

\begin{frame}
	Jekyll uses Sass as the main filetype for defining the styles to be applied to our HTML. These Sass files are processed into a CSS file during the Jekyll `build phase'.
	
	\centering
	\includegraphics[width=0.5\textwidth]{jekyllvial}
\end{frame}

% subsection sass (end)

\subsection*{Javascript and Coffeescript} % (fold)
\label{sub:javascript_and_coffeescript}

\begin{frame}{Client-Side Scripts}
	\begin{itemize}
		\item Javascript (or Coffeescript -- processed to Javascript) is the client-side programming language of the web.
		\item JavaScript can be used to manipulate the content, layout and styling of a web page based on the users actions.
	\end{itemize} 
\end{frame}
% subsection javascript_and_coffeescript (end)

\subsection*{Summary} % (fold)
\label{sub:summary}

\begin{frame}{Summary}
	\begin{itemize}
		\item<1-> We write our content using Markdown which is processed by Jekyll to produce the HTML.
		\item<2-> We write our styling using SASS which is processed into CSS.
		\item<3-> In practice, we can use a theme which defines the style for us. 
		\item<3-> We can write our client-side scripts as JavaScript (or CoffeeScript).
	\end{itemize}
\end{frame}
% subsection summary (end)

% section bblocks (end)

\section{Jekyll Demo} % (fold)
\label{sec:jekyll}

\begin{frame}{Building a site with Jekyll}
	\Huge \centering Demo
	
	\medskip
	\normalsize
	This demo uses the \href{https://github.com/gaalcaras/academic}{Academic} theme by Gabriel Alcaras.
\end{frame}
% section jekyll (end)

\end{document}